\subsubsection{Regional continuation and renewal equation}

The selector \(169\) determines the degree-six impulse. Subsequent
prolongation is organized by the transfer \(D_A\) and a unit
normalization of the determinant line. These rules admit a formulation
in formal series before substituting \(z=\alpha\), with \(D_A\) already
constructed. Write
\[
 \mathscr C(z)=169z^6+R(z),\qquad R(z)\in z^7\mathbb R[[z]].
\]
Concatenating a prolongation increases the degree by one unit and
multiplies its weight by \(D_A\). The first strict prolongation has
weight \(D_A\). Its renewal equation is therefore
\begin{equation}
 R(z)=D_Az^7+D_AzR(z).
 \label{eq:revision-renovacion}
\end{equation}

\begin{proposition}[Formal uniqueness of the normalized continuation]
Equation \eqref{eq:revision-renovacion} determines a unique series, and
\[
 \mathscr C(z)=169z^6+\frac{D_Az^7}{1-D_Az}
             =169z^6+\sum_{n=7}^{\infty}D_A^{n-6}z^n.
\]
The realization is analytic in \(|D_Az|<1\).
\end{proposition}
\begin{proof}
The element \(1-D_Az\) has constant term one and is invertible
in the ring of formal series. Solving the equation produces the displayed
series. Equivalently, the coefficient of degree seven is \(D_A\),
and each subsequent coefficient is \(D_A\) times the preceding one. The geometric
series converges absolutely in the indicated disk.
\end{proof}

Normalization of the first prolongation is an indispensable
operational datum. If only the degree-six impulse and
concatenation ratio were retained, the series
\[
 169z^6+\lambda\frac{D_Az^7}{1-D_Az}
\]
would continue to share both data for every \(\lambda\).
The unit rule fixes \(\lambda=1\) before evaluation.
The contribution of each rule to uniqueness of the reader is thus
localized, and the rational expression preserves the complete continuation.

\subsubsection{Reduced Planck constant and return of the action section}

In the dimensional chart of \eqref{eq:el-secciones-accion}, the
section before return is the model construction intended for comparison
with the reduced Planck constant:
\[
 \hbar_{\mathrm{pre}}=H_5\,10^{-34}\mathcal U_S.
\]
The subsequent section additionally preserves the regional compensation
\(90\mathscr C(\alpha)/\pi\). Both belong to the same action
line, and their ratio \(R_{\mathrm{act}}\) measures the multiplicative effect
of return. Passage to \(h\) through \(2\pi\) uses the relation
between angular parameterization and a complete turn.

This construction effectively enters
\(R_{\mathrm{act}}^{80-54\alpha+6\Delta_4}\). The electron preserves
its central coordinate and trajectory; the chronological scale does not
replace that structure with an undifferentiated quantum. Nor is a
mass already expressed in MeV multiplied again by \(10^{-34}\):
the decade belongs to the action section, cancels in its dimensionless
ratio, and reappears only where the corresponding dimensional
realization acts.

The numerical comparison of \(H_5\), the returned section, and
\(h/(2\pi)\) is presented in Section~\ref{sec:revision-planck-contraste}.
This distinguishes the name of the physical object, the equation the
model assigns to it, and the precision with which that equation reproduces it.
